\begin{textsample}{POS Dim 2 – pb – Score 28.00 – pb/pb\_em\_2.txt}  \label{ex:f2_pos_261}
1.1 APRESENTAÇÃO A construção de uma Proposta Curricular do Ensino Médio para o Estado da Paraíba apresenta -se como uma ação necessária, \textbf{atendendo} ao \textbf{disposto} no Art. 210 da Constituição Federal ( 1988 ) sobre a definição de uma Base Nacional Comum com conteúdos mínimos, de maneira a assegurar formação básica comum.
Nessa linha, a Lei de \textbf{Diretrizes} e Bases da Educação Nacional ( LDB ) de nº 9.394, de 20 de dezembro de 1996, no Art. 26 retoma a necessidade de uma Base Comum Curricular a ser complementada em cada sistema de ensino e em cada estabelecimento escolar.
Assim como também de forma a seguir a propositura da Lei n{\EmojiFont °} 13.005/2014, Plano Nacional de Educação ( PNE ), Meta 7, no que diz respeito a fomentar a qualidade da educação básica e da Lei nº 10.488, de 23 de junho de 2015, Plano Estadual de Educação ( PEE ) em sua Meta 4, corresponde à Meta 3 do PNE, Estratégia 4.2 no que preconiza contribuir com o MEC na elaboração da proposta de direitos e objetivos de aprendizagem e desenvolvimento para os ( as ) alunos ( as ) de ensino médio a serem atingidos nos tempos e etapas de organização desta etapa de ensino, com vistas a garantir formação básica comum, incluindo a participação democrática da sociedade civil.
Para a formação básica comum e parte diversificada, as \textbf{Diretrizes} Curriculares Nacionais ( DCNEN ) define a Base Nacional Comum, no seu Art. 14, “ como conhecimentos, saberes e valores produzidos culturalmente, [... ] que são gerados nas \textbf{instituições} produtoras do conhecimento científico e tecnológico ; no mundo do trabalho ; no desenvolvimento das linguagens ; nas atividades desportivas e corporais ; na produção artística ; na cidadania ”.
Com a homologação da Base Nacional Comum Curricular ( BNCC ), a educação toma um novo rumo, isto é, \textbf{normatiza} a sistematização das aprendizagens essenciais que todos os estudantes devem desenvolver ao longo da educação básica – de forma progressiva e por áreas de conhecimento, enquanto que os \textbf{currículos} traçam os caminhos e são implementados nos sistemas de ensino, considerando o \textbf{currículo} local, às suas características regionais, diante da diversidade histórica, social, ambiental e cultural que exigem da Educação Básica novas concepções e arranjos curriculares capazes de dialogar com as juventudes e seus anseios.
Nesse pressuposto, a Proposta Curricular do Ensino Médio do Estado da Paraíba é resultado de um trabalho colaborativo, que envolveu a consulta aos estudantes do 9º ano do Ensino Fundamental e das três \textbf{séries} do Ensino Médio, professores e gestores escolares das diferentes modalidades \textbf{ofertadas} na rede estadual, por meio de seminários de escuta e leituras críticas aos documentos curriculares \textbf{preliminares} disponibilizados nas consultas públicas.
O trabalho de redação resultou da construção realizada pelos professores da rede estadual de ensino e suas experiências e atuações em escolas com \textbf{ofertas} em diferentes modalidades no Ensino Médio : Regular, Integral, Educação Profissional e Tecnológica, Educação de Jovens e Adultos, Educação Especial, Educação do Campo, Indígena, Cigana e Quilombola.
Os professores \textbf{redatores} foram coordenados por uma \textbf{equipe} técnico-pedagógica composta por profissionais da Secretaria de Estado da Educação e da Ciência e Tecnologia da Paraíba - SEECT/PB com apoio \textbf{técnico} do Conselho Nacional de \textbf{Secretários} de Educação - CONSED e União Nacional dos \textbf{Dirigentes} \textbf{Municipais} de Educação - Paraíba - UNDIME/PB e Conselho Estadual da Educação da Paraíba - CEE/PB.
A integração de diferentes atores permitiu a que o processo de elaboração da Proposta Curricular do Ensino Médio do Estado da Paraíba não estivesse centralizada em apenas um ator.
Por meio de consultas voluntárias realizadas às \textbf{Instituições} Públicas de Ensino Superior da Paraíba e parceiros, com apoio do CONSED, foi possível concluir um documento plural, com contribuições que partem de diferentes perspectivas e leituras críticas.
Foram dois anos de trabalho ( 2019/2020 ) com formações envolvendo palestras, oficinas, práticas de escritas, escutas qualificadas, revisões e reescritas para a composição do formato e estrutura atual do documento que agora apresentamos, sem a pretensão de ser um documento fechado, pelo contrário com ponto de partida para novos estudos, proposições e atualizações ao entender o \textbf{currículo} como um documento que expresse as transformações de seu tempo.
Ao longo desta construção foi possível sugerir um arranjo pedagógico dentro do que disciplina a \textbf{legislação} educacional em vigor, contemplando as especificidades do território da Paraíba, de modo a proporcionar aos estudantes a possibilidade de vivências diversas a partir de seus interesses e fortalecendo seus projetos de vida.
Aos professores, possibilidades de reflexão sobre seu campo de atuação e sobre a condução pedagógica referente a etapa do Ensino Médio em constante e necessário diálogo com \textbf{gestão} escolar, apontando caminhos da organização de \textbf{oferta} curricular com a Formação Geral Básica, que propõe a aprendizagem das competências e habilidades definidas pela Base Nacional Comum Curricular ( BNCC ), e dos \textbf{Itinerários} Formativos, que compreende um conjunto de unidades curriculares que os estudantes podem escolher a partir do seu interesse para aprofundar e ampliar aprendizagens em uma ou mais Áreas de Conhecimento e/ou na Formação \textbf{Técnica} e Profissional, promovendo o \textbf{fortalecimento} da construção e/ou atualização dos Projetos Políticos Pedagógicos ( PPPs ) e do planejamento do professor de cada unidade escolar do território paraibano.
A elaboração da Proposta Curricular do Ensino Médio da Paraíba ocorreu a partir de um processo coletivo no qual muitas mãos se estabeleceram como coautoras das páginas que seguem.
Não corresponde a um processo mecânico ou burocrático, mas apresenta -se como a coleção de vários momentos de reflexão, angústias, indagações, expectativas, estudos, de muito trabalho.
Cada \textbf{redator} trouxe para seu texto sua vivência enquanto professor da Educação Básica, suas escutas cotidianas das vozes de cada estudante que formam, entrelaçando com a necessidade de responder a uma pergunta essencial : Qual escola queremos para os \textbf{adolescentes}, jovens e adultos paraibanos?
Como a Proposta Curricular do Ensino Médio da Paraíba, no contexto da escola, pode ajudar na formação do cidadão que queremos?

% matched lemmas: adolescente, atender, currículo, diretriz, dirigente, disposto, equipe, fortalecimento, gestão, instituição, itinerário, legislação, municipal, normatizar, oferta, ofertar, preliminar, redator, secretário, série, técnico
\end{textsample}
